\begin{textsample}{POS Dim 2 – mg – Score 35.00 – mg/mg\_em\_103.txt}  \label{ex:f2_pos_191}
5.1 Apresentação A educação profissional e tecnológica é uma modalidade de ensino que busca preparar o estudante para o exercício do trabalho e sua inserção no mercado.
As grandes mudanças nessa \textbf{oferta} estão presentes de forma mais intensa no Brasil a partir da metade do século XX em diante, quando a indústria se potencializa e se torna um forte pilar para a economia brasileira.
A primeira Lei de \textbf{Diretrizes} e Base da Educação Nacional, Lei nº 4.024, de 20 de dezembro de 1961, foi debatida e elaborada no contexto de redemocratização do país, porém só foi promulgada como Lei nº 4.024 em 20 de dezembro de 1961.
O sistema de ensino continuou a ser organizado segundo as \textbf{legislações} anteriores, mantendo a divisão dos \textbf{cursos} secundário, \textbf{técnico} e \textbf{curso} normal.
A partir da promulgação da Lei nº 5.692, de 11 de agosto de 1971, o ensino técnico-profissional passou a integrar, de forma compulsória, o \textbf{currículo} do segundo grau.
Essa \textbf{reforma} de ensino decorreu da \textbf{política} governamental voltada para a contenção da demanda do ensino superior, considerada incompatível com as necessidades nacionais.
A formação \textbf{técnica} obrigatória impulsionou o número de matrículas nas redes públicas e privadas, e também a implantação de novos \textbf{cursos} \textbf{técnicos}.
As disposições legais da profissionalização do ensino de 2º Grau foram disciplinadas por meio do \textbf{Parecer} nº 45/1972 e da Resolução nº 08/1972, do Conselho Federal de Educação, que estabeleceram o \textbf{currículo} mínimo e classificaram a \textbf{qualificação} para o trabalho como um componente básico do processo de formação integral do estudante nos ensinos de 1º e 2º graus.
Nesse processo ficou determinado que, no Ensino de 1º grau, ela tomasse a forma de sondagem de aptidões e iniciação para o trabalho, enquanto no Ensino de 2º Grau tivesse papel preponderante para a \textbf{habilitação} profissional.
Deveria ser oferecida em \textbf{cursos} profissionalizantes de 2º Grau com duração de três ou quatro \textbf{séries}.
No entanto, as dificuldades materiais de implantação do ensino profissionalizante nas escolas de segundo grau, tanto públicas quanto particulares, somadas às resistências de vários segmentos, fizeram com que a Lei nº 5.692/1971 fosse reinterpretada pelo \textbf{Parecer} nº 76/1975, do Conselho Federal de Educação, que manteve como objetivo do segundo grau o ensino profissionalizante, mas mudou seu conteúdo.
A \textbf{habilitação} profissional passou a chamar -se específica, isto é, passou a voltar -se para diferentes especialidades \textbf{técnicas}, sendo entendida como meio de tornar o jovem consciente do domínio que deveria ter das bases científicas de uma profissão, além de se tornar apto à aplicação da tecnologia adequada.
A Lei nº 7.044, de 23 de agosto de 1982, alterou dispositivos da Lei 5.692/1971, resgatando a ênfase da formação geral e dispensando a obrigatoriedade da profissionalização.
Assim, a formação profissional, em nível de segundo grau, ficou a cargo de escolas \textbf{técnicas}, tal como antes, sendo oferecida com \textbf{cursos} longos ou compactos para aqueles que concluíssem alguma \textbf{habilitação} básica ou, ainda, por treinamento em serviço.
Nos anos 1990, o Brasil foi marcado pela mobilização social por \textbf{reformas} no setor educacional que \textbf{culminaram} com a elaboração da atual Lei de \textbf{Diretrizes} e Bases da Educação Nacional, Lei nº 9.394, de 20 de \textbf{novembro} de 1996.
A partir de então, o segundo grau é \textbf{instituído} como ensino médio e são incluídas disposições específicas para o desenvolvimento da Educação Profissional em um \textbf{capítulo} à parte da Educação Básica.
Acompanhando os princípios da Constituição Federal de 1988, a LDB 9.394/1996 reconhece que a educação abrange os processos formativos que se desenvolvem nas distintas esferas de convivência social e esclarece que a \textbf{legislação} disciplina apenas a educação escolar, desenvolvida em \textbf{instituições} próprias, e que a educação deve vincular -se ao mundo do trabalho e à prática social.
As finalidades da educação nacional são reafirmadas no \textbf{artigo} 2º, isto é, o pleno desenvolvimento do educando, seu preparo para o exercício da cidadania e sua \textbf{qualificação} para o trabalho.
Nessa perspectiva, articula formação geral e profissional, de modo \textbf{flexível} – podendo ser obtida a formação profissional a partir de diferentes níveis e \textbf{instituições} de ensino.
Em 2008, a Lei nº 11.741, de 16 de \textbf{julho} de 2008, alterou dispositivos da Lei no 9.394/1996 quanto à educação profissional.
Destaca -se : - A integração da educação profissional e tecnológica aos diferentes níveis e modalidades de educação e às dimensões do trabalho, da ciência e da tecnologia. - A organização dos \textbf{cursos} de educação profissional e tecnológica por eixos tecnológicos, possibilitando a construção de diferentes \textbf{itinerários} formativos. - A configuração e abrangência dos \textbf{cursos} de educação profissional e tecnológica : I ) formação inicial e continuada ou \textbf{qualificação} profissional ; II ) \textbf{técnica} de nível médio ; III ) tecnológica de \textbf{graduação} e pós-graduação. - O conhecimento adquirido na educação profissional e tecnológica, inclusive no trabalho, poderá ser objeto de avaliação, reconhecimento e \textbf{certificação} para \textbf{prosseguimento} ou conclusão de estudos. - A preparação geral para o trabalho e, facultativamente, a \textbf{habilitação} profissional, poderão ser desenvolvidas nos próprios estabelecimentos de ensino médio ou em cooperação com \textbf{instituições} especializadas em educação profissional. - A educação profissional \textbf{técnica} de nível médio será desenvolvida nas formas articuladas com o ensino médio ( integrado ou concomitante ) ou para quem já tenha concluído o ensino médio ( subsequente ). - Os diplomas de \textbf{cursos} de educação profissional \textbf{técnica} de nível médio, quando registrados, têm validade nacional e habilitam ao \textbf{prosseguimento} de estudos na educação superior. - Os \textbf{cursos} \textbf{técnicos}, nas formas articulada, concomitante e subsequente, quando estruturados e organizados em etapas com terminalidade possibilitarão a obtenção de certificados de \textbf{qualificação} para o trabalho após a conclusão ( \textbf{certificações} intermediárias ).
Em 2006, o Decreto nº 5.840, de 13 de \textbf{julho} de 2006, \textbf{institui}, no âmbito \textbf{federal}, o Programa Nacional de Integração da Educação Profissional com a Educação Básica na Modalidade – PROEJA, abrangendo o ensino fundamental, médio e a educação indígena.
A fim de regulamentar a \textbf{oferta} de \textbf{cursos} \textbf{técnicos} de nível médio, o \textbf{Catálogo} Nacional dos \textbf{Cursos} \textbf{Técnicos} - CNCT, elaborado durante o ano de 2007, esteve em consulta pública por seis meses. \textbf{Instituído} pela \textbf{Portaria} MEC nº 870, de 16 de \textbf{julho} de 2008, com base no \textbf{Parecer} CNE/CEB nº 11/2008 e na Resolução CNE/CEB nº 3/2008, é \textbf{atualizado} periodicamente para contemplar novas demandas socioeducacionais.
A quarta versão do CNCT está \textbf{prevista} para consulta e publicação em 2020.
Outro marco legal importante para a Educação Profissional no país foi a Resolução nº 06, de 20 de setembro de 2012, que definiu as \textbf{Diretrizes} Curriculares Nacionais para a Educação Profissional \textbf{Técnica} de Nível Médio.
Nessa normativa estão definidos os princípios norteadores da Educação Profissional \textbf{Técnica} de Nível Médio ; a organização curricular e formas de \textbf{oferta} ; duração dos \textbf{cursos} ; avaliação, aproveitamento e \textbf{certificação} ; e, ainda, \textbf{diretrizes} para a formação docente.
Assim como o CNCT, a Resolução nº 06/2012 foi revista e \textbf{atualizada} pela Secretaria de Educação Profissional e Tecnológica – SETEC do Ministério da Educação em 2020.
A Resolução CNE/CP nº 1, de 05 de \textbf{janeiro} de 2021, \textbf{atualizou} e definiu as \textbf{Diretrizes} Curriculares Nacionais Gerais para a Educação Profissional e Tecnológica.
Paralelamente aos avanços legais da Educação Profissional e Tecnológica no âmbito da Educação Nacional, o Ministério do Trabalho e Emprego implementou nas últimas décadas uma \textbf{série} de propostas que alcançam \textbf{instituições} e agentes da Educação Profissional de todo o país ( universidades, redes de ensino \textbf{técnico} público e privado, \textbf{instituições} de formação profissional, escolas livres, sindicatos, empresas, ONGs ).
Dentre eles estão o Plano Nacional de Educação Profissional – PLANFOR ; o Fundo de Amparo ao \textbf{Trabalhador} – FAT, os Planos Estaduais de \textbf{Qualificação} – PEQ ; protocolos e acordos de cooperação \textbf{técnica} e financeira ; avanço metodológico, conceitual e operacional em ações de \textbf{qualificação} profissional.

% matched lemmas: artigo, atualizar, capítulo, catálogo, certificação, culminar, currículo, curso, diretriz, federal, flexível, graduação, habilitação, instituir, instituição, itinerário, janeiro, julho, legislação, novembro, oferta, parecer, política, portaria, prever, prosseguimento, qualificação, reforma, série, trabalhador, técnico
\end{textsample}
